\documentclass[11pt]{amsart}
\usepackage[T1]{fontenc}
\usepackage{lmodern}
\usepackage{amsmath,amssymb,amsthm,mathtools}
\usepackage{xcolor}
\usepackage[colorlinks=true,linkcolor=blue!55!black,citecolor=blue!55!black,urlcolor=blue!55!black]{hyperref}

\newtheorem{theorem}{Theorem}[section]
\newtheorem{lemma}[theorem]{Lemma}
\newtheorem{proposition}[theorem]{Proposition}
\newtheorem{corollary}[theorem]{Corollary}
\theoremstyle{definition}
\newtheorem{definition}[theorem]{Definition}
\newtheorem{remark}[theorem]{Remark}

\title{A Counterexample to an Exponential Divisibility Recurrence}
\author{Xiaojian Zeng}
\address{PBC School of Finance, Tsinghua University, 100190}
\email{xiaojianzeng@foxmail.com}
\date{2026-09-02}
\hypersetup{pdftitle={A Counterexample to an Exponential Divisibility Recurrence},pdfauthor={Xiaojian Zeng}}

\begin{document}

\begin{abstract}
Starting with $a_1=a_2=1$, suppose that
$a_n=(2^{a_{n-1}}+1)/a_{n-2}$ whenever this quotient is integral.
We show that all divisions through $a_6$ are exact but the next one is not:
$(2^{a_6}+1)\bmod 171=9$.  Thus the integrality question recorded as
OEIS A078590 has the negative answer, with first failure at $n=7$.
\end{abstract}

\maketitle

\section{Introduction}
The entry OEIS A078590, contributed by Benoit Cloitre in 2002, asks whether
the recurrence
\begin{equation}\label{eq:rec}
 a_1=a_2=1,\qquad a_n=\frac{2^{a_{n-1}}+1}{a_{n-2}}\quad(n\geq3)
\end{equation}
always produces integers \cite{oeis}.  The public entry displays the terms
through $a_6$ and still poses the integrality question.  The corresponding
Formal Conjectures source also records the divisibility assertion as open
\cite{formalconjectures}.  We give an elementary exact counterexample at the
very next step.  Our literature statement is limited to those public records
and searches completed on September 2, 2026; it is not a claim that no
unindexed or differently worded prior observation exists.

Equation \eqref{eq:rec} is understood as a partial recurrence: $a_n$ is
defined only if all preceding terms are positive integers and the displayed
division is exact.  This matters because totalizing division in the natural
numbers would define a different sequence after the first failure.

\section{Main result}

\begin{theorem}\label{thm:main}
For the recurrence \eqref{eq:rec}, the terms $a_1,\ldots,a_6$ are positive
integers, but
\[
 171\nmid 2^{a_6}+1.
\]
More precisely,
\[
 2^{a_6}+1\equiv9\pmod{171}.
\]
Consequently the first failed division is the one defining $a_7$.
\end{theorem}

\begin{proof}
Direct calculation gives
\[
 a_3=3,\qquad a_4=9,\qquad a_5=171.
\]
As $2^3\equiv-1\pmod9$ and $171=3\cdot57$, with $57$ odd, we have
$2^{171}\equiv-1\pmod9$.  Hence $a_6=(2^{171}+1)/9$ is an integer.

It remains to test the next numerator without expanding it.  First,
\[
 2^{171}+1\equiv27\pmod{162}.
\]
(For example, modulo $81$ the exponent may be reduced modulo $54$, giving
$2^{171}\equiv2^9\equiv26\pmod{81}$; parity then selects the residue $27$
modulo $162$ after adding one.)  Division of this congruence by $9$ is now
legitimate and yields
\[
 a_6=\frac{2^{171}+1}{9}\equiv3\pmod{18}.
\]
Thus $a_6\equiv3\pmod6$, so $2^{a_6}+1\equiv2^3+1\equiv0\pmod9$.
Also Fermat's theorem modulo $19$, together with $a_6\equiv3\pmod{18}$,
gives $2^{a_6}+1\equiv9\pmod{19}$.  Since $171=9\cdot19$, the Chinese
remainder theorem gives
\[
 2^{a_6}+1\equiv9\pmod{171}.
\]
The nonzero remainder proves failure at $n=7$, while the preceding
calculations prove that no earlier division fails.
\end{proof}

\section{Independent exact verification}

For reproducibility, the ancillary files contain the integer $a_6$, the
remainder, and two small verifiers.  One reconstructs all terms through
$a_6$ using exact integer arithmetic; the other implements binary modular
exponentiation independently.  Both return the exact remainder $9$.  These
programs are supplementary checks: the proof above is complete without them.

\section*{Data, code, and computational disclosure}
The certificate and verifiers accompanying this note use Python 3.10 and
Ruby 3.0 and require no external solver or floating-point computation.
Automated language assistance was used during drafting and adversarial
checking; the exact argument and all scientific claims were independently
recomputed and are the author's responsibility.

\bibliographystyle{amsplain}
\bibliography{references}

\end{document}
